\setcounter{section}{22}

  \zwischenueberschrift{Übungsaufgaben}

\inputaufgabe
{}
{

Es sei $M$ eine
\definitionsverweis {berandete Mannigfaltigkeit}{}{}
und
\mathl{\partial M}{} sei der Rand. Zeige, dass der
\definitionsverweis {topologische Rand}{}{}
von
\mathl{M \setminus \partial M}{} gleich $\partial M$ ist.

}
{} {}

\inputaufgabe
{}
{

Bestimme die
\definitionsverweis {Träger}{}{} der folgenden Funktionen von $\R$ nach $\R$.
\aufzaehlungsieben{Eine Polynomfunktion.
}{Die Sinusfunktion.
}{Die Exponentialfunktion.
}{Die Indikatorfunktion $e_{ \Z }$.
}{Die Indikatorfunktion $e_{ \Q }$.
}{Die Indikatorfunktion $e_{  [a,b]  }$.
}{Die Indikatorfunktion $e_{ ]a,b[ }$.
}

}
{} {}

\inputaufgabe
{}
{

Es sei
\maabbdisp {f} { X } {\R
} {}
eine Funktion auf einem
\definitionsverweis {topologischen Raum}{}{}
$X$. Es gebe eine
\definitionsverweis {offene Teilmenge}{}{}

\mavergleichskette
{\vergleichskette
{U
}
{ \subseteq }{ X
}
{  }{
}
{    }{
}
{   }{
}
}
{}{}{,}
die den
\definitionsverweis {Träger}{}{}
von $f$ umfasse. Zeige, dass $f$ genau dann
\definitionsverweis {stetig}{}{}
ist, wenn die Einschränkung
\mathl{f \vert_U}{} stetig ist.

}
{} {}

\inputaufgabe
{}
{

Es sei $X$ ein
\definitionsverweis {topologischer Raum}{}{}
und

\mavergleichskette
{\vergleichskette
{ T
}
{ \subseteq }{ X
}
{  }{
}
{    }{
}
{   }{
}
}
{}{}{}
eine Teilmenge. Zeige, dass der
\definitionsverweis {Abschluss}{}{}
von $T$ gleich dem
\definitionsverweis {Träger}{}{}
der
\definitionsverweis {Indikatorfunktion}{}{}

\mathl{e_{  T  }}{} ist.

}
{} {}

\inputaufgabe
{}
{

Man gebe eine
\definitionsverweis {kompakte Ausschöpfung}{}{} für die
\definitionsverweis {reellen Zahlen}{}{} $\R$ an.

}
{} {}

\inputaufgabegibtloesung
{}
{

Man gebe eine
\definitionsverweis {kompakte Ausschöpfung}{}{}
für den $\R^d$ an.

}
{} {}

\inputaufgabe
{}
{

Es sei $X$ ein
\definitionsverweis {topologischer Raum}{}{.}
Bestimme zur Überdeckung von $X$ durch $X$ eine
\definitionsverweis {untergeordnete Partition der Eins}{}{.}

}
{} {}

\inputaufgabe
{}
{

Es sei $X$ ein
\definitionsverweis {topologischer Raum}{}{}
und

\mavergleichskette
{\vergleichskette
{ X
}
{ = }{ \bigcup_{i \in I} U_i
}
{  }{
}
{    }{
}
{   }{
}
}
{}{}{}
eine
\definitionsverweis {offene Überdeckung}{}{.}
Wir betrachten die Familie der
\definitionsverweis {Indikatorfunktionen}{}{}

\mathbeddisp {e_{  P  }} {}
 {P \in X} {}
 {} {} {} {.}
Welche Eigenschaften einer
\zusatzklammer {dieser Überdeckung} {} {}
\definitionsverweis {untergeordneten Partition der Eins}{}{}
erfüllt diese Familie?

}
{} {}

\inputaufgabe
{}
{

Wir betrachten die
\definitionsverweis {kompakte Ausschöpfung}{}{}

\mathbed {A_n= [-n,n]} {}
 {n \in \N} {}
 {} {} {} {,} der
\definitionsverweis {reellen Zahlen}{}{}
und die offene Überdeckung

\mathbed {W_n= A_{n+1}^{o}  \setminus A_{n-1}} {}
 {n \in \N} {}
 {} {} {} {,}
\zusatzklammer {es sei
\mathl{A_{-1}= \emptyset}{}} {} {.}
Finde eine Überdeckung von $\R$ mit
\definitionsverweis {offenen Intervallen}{}{,}
die die Eigenschaften
aus [[Mannigfaltigkeit/Abzählbare Basis/Lokal endlicher Atlas/Fakt|Lemma 22.9]]
\zusatzklammer {und seinem Beweis} {} {} erfüllt.

}
{} {}

\inputaufgabe
{}
{

Man konstruiere eine Folge von
\definitionsverweis {stetigen Funktionen}{}{}
\maabbdisp {f_n} { [0,1] } { [0,1]
} {}
zu

\mavergleichskette
{\vergleichskette
{ n
}
{ \in }{ \N
}
{  }{
}
{    }{
}
{   }{
}
}
{}{}{,}
die eine
\definitionsverweis {Partition der Eins}{}{}
bilden und die Eigenschaft erfüllen, dass $0$ zum
\definitionsverweis {Träger}{}{}
einer jeden Funktion $f_n$ gehört.

}
{} {}

\inputaufgabe
{}
{

Es sei

\mavergleichskette
{\vergleichskette
{ V
}
{ \subseteq }{ H
}
{  }{
}
{    }{
}
{   }{
}
}
{}{}{}
eine
\definitionsverweis {offene Menge}{}{}
im
\definitionsverweis {Halbraum}{}{}

\mavergleichskette
{\vergleichskette
{ H
}
{ \subset }{ \R^n
}
{  }{
}
{    }{
}
{   }{
}
}
{}{}{}
und sei
\maabbdisp {f} { V } {\R
} {}
eine
\definitionsverweis {stetig differenzierbare Funktion}{}{.}
Zeige, dass es eine offene Menge

\mavergleichskette
{\vergleichskette
{ \tilde{V}
}
{ \subseteq }{ \R^n
}
{  }{
}
{    }{
}
{   }{
}
}
{}{}{}
und eine
\definitionsverweis {stetig differenzierbare Funktion}{}{}
\maabbdisp {\tilde{f}} {\tilde{V}} { \R
} {}
mit

\mavergleichskette
{\vergleichskette
{ \tilde{V} \cap H
}
{ = }{ V
}
{  }{
}
{    }{
}
{   }{
}
}
{}{}{}
und

\mavergleichskette
{\vergleichskette
{ \tilde{f} \vert_{\tilde{V} }
}
{ = }{ f
}
{  }{
}
{    }{
}
{   }{
}
}
{}{}{}
gibt.

}
{} {}

\inputaufgabe
{}
{

Es sei
\maabbdisp {f} {\R} {\R
} {}
eine
\definitionsverweis {stetig differenzierbare Funktion}{}{}
mit
\definitionsverweis {kompaktem Träger}{}{.}
Zeige, dass die
\definitionsverweis {Ableitung}{}{}
$f'$ ebenfalls kompakten Träger hat, und dass

\mavergleichskette
{\vergleichskette
{ \int_\R f' d \lambda^1
}
{ = }{ 0
}
{  }{
}
{    }{
}
{   }{
}
}
{}{}{}
ist.

}
{} {}

\inputaufgabe
{}
{

Es sei
\maabbdisp {f} {\R_{\geq 0}} {\R
} {}
eine
\definitionsverweis {stetig differenzierbare Funktion}{}{}
mit
\definitionsverweis {kompaktem Träger}{}{.}
Zeige, dass die
\definitionsverweis {Ableitung}{}{}
$f'$ ebenfalls kompakten Träger hat, und dass

\mavergleichskette
{\vergleichskette
{ \int_{\R_{\geq 0} } f' d \lambda^1
}
{ = }{ -f(0)
}
{  }{
}
{    }{
}
{   }{
}
}
{}{}{}
ist. Zeige, dass diese Aussage nicht gelten muss, wenn $f$ keinen kompakten Träger besitzt.

}
{} {}

\inputaufgabe
{}
{

Es sei $X$ eine
\definitionsverweis {differenzierbare Mannigfaltigkeit}{}{}
mit einer
\definitionsverweis {abzählbaren Basis der Topologie}{}{.}
Zeige, dass es eine
\definitionsverweis {riemannsche Struktur}{}{}
auf $X$ gibt.

}
{} {}

\inputaufgabe
{}
{

Es sei $X$ eine
\definitionsverweis {differenzierbare Mannigfaltigkeit}{}{}
mit einer
\definitionsverweis {abzählbaren Basis der Topologie}{}{}
und sei
\maabb {p} { E } { X
} {}
ein
\definitionsverweis {differenzierbares Vektorbündel}{}{}
über $X$. Zeige, dass man $E$ mit der Struktur eines
\definitionsverweis {riemannschen Vektorbündels}{}{}
versehen kann.

}
{} {}

  \zwischenueberschrift{Aufgaben zum Abgeben}

\inputaufgabe
{3}
{

Es sei

\mathbed {A_n} {}
 {n \in \N} {}
 {} {} {} {,}
eine
\definitionsverweis {kompakte Ausschöpfung}{}{} eines
\definitionsverweis {topologischen Raumes}{}{} $X$. Zeige, dass die Beziehung

\mathdisp {A_{n+1} \setminus A_n ^{o}  \subseteq A_{n+2}^{o} \setminus A_{n-1}} {  }
 gilt.

}
{} {}

\inputaufgabe
{4}
{

Man gebe zur
\definitionsverweis {offenen Überdeckung}{}{}

\mavergleichskettedisp
{\vergleichskette
{ \R
}
{ =} { \bigcup_{n \in \Z} ]n,n+3[
}
{ } {
}
{ } {
}
{ } {
}
}
{}{}{}
eine untergeordnete stetige
\definitionsverweis {Partition der Eins}{}{}
an.

}
{} {}

\inputaufgabe
{6}
{

Wir betrachten die
\definitionsverweis {kompakte Ausschöpfung}{}{}

\mathbeddisp {A_n = B  \left(  0 ,n \right)} {}
 {n \in \N} {}
 {} {} {} {,}
des $\R^2$ und die
\definitionsverweis {offene Überdeckung}{}{}

\mathbeddisp {W_n= A_{n+1}^{o} \setminus A_{n-1}} {}
 {n \in \N} {}
 {} {} {} {,}
\zusatzklammer {es sei
\mathl{A_{-1}= \emptyset}{}} {} {.}
Finde eine Überdeckung des $\R^2$ mit
\definitionsverweis {offenen Kreisscheiben}{}{,}
die die Eigenschaften
aus [[Mannigfaltigkeit/Abzählbare Basis/Lokal endlicher Atlas/Fakt|Lemma 22.9]]
\zusatzklammer {und seinem Beweis} {} {}
erfüllt.

}
{} {}

\inputaufgabe
{3}
{

Man gebe ein Beispiel für einen
\definitionsverweis {topologischen Raum}{}{,}
der keine
\definitionsverweis {kompakte Ausschöpfung}{}{}
besitzt.

}
{} {}

 [[Kategorie:Latexseite]]
